% =====================================================================
% Chương 2 — Lớp Mạng (Network)
% Phân tích: espnow_protocol.h, espnow_client.cpp, networkTask (main.cpp),
%            espnow_receiver.c, coreiot_client (2 board), platformio.ini
% =====================================================================

\chapter{Lớp Mạng (Network)}

Lớp Mạng vận chuyển dữ liệu đã lọc từ Lớp 1 tới đích bằng \textbf{hai tuyến
song song}: ESP-NOW (đường chính trên xe) và MQTT/Wi-Fi (đường phụ lên cloud).
Cả hai tuyến đều đọc cùng một nguồn \code{shared\_state} nhưng gói tin và tần
suất khác nhau.

\section{Tổng quan hai tuyến}

\begin{center}
\small
\renewcommand{\arraystretch}{1.15}
\begin{longtable}{@{}p{4.0cm}p{5.6cm}p{6.0cm}@{}}
\caption{So sánh đường chính và đường phụ.}\\
\toprule
\textbf{Tiêu chí} & \textbf{ESP-NOW (chính)} & \textbf{MQTT (phụ)} \\
\midrule
\endhead
Giao thức & Link Layer Espressif, không AP & TCP/MQTT qua Wi-Fi \\
Tần suất & \code{ESPNOW\_SEND\_INTERVAL\_MS 500} & \code{COREIOT\_PUBLISH\_INTERVAL\_MS 500} \\
Payload & \code{espnow\_sensor\_msg\_t} packed 30\,B & JSON \code{d1..d6 + nearest\_cm} \\
Độ trễ mục tiêu & <50\,ms & 100--500\,ms \\
Phụ thuộc Internet & Không & Có \\
Đích & Waveshare-screen & app.coreiot.io:1883 \\
\bottomrule
\end{longtable}
\end{center}

\section{Hợp đồng dây: espnow\_protocol.h (R2)}

File duy nhất định nghĩa giao diện ESP-NOW; cấm định nghĩa lại symbol ở nơi
khác (R2 — \code{grep} toàn firmware = 1). Các hằng số quan trọng:

\begin{verbatim}
#define ESPNOW_CHANNEL 1
static const uint8_t ESPNOW_PEER_MAC[6] = {0x64,0xe8,0x33,0x7c,0x3f,0xe0};
#define ESPNOW_SEND_INTERVAL_MS 500
#define ESPNOW_LINK_TIMEOUT_MS 1500
typedef enum { ESPNOW_SLOT_FRONT=0, ESPNOW_SLOT_REAR=1,
               ESPNOW_SLOT_LEFT_FRONT=2, ESPNOW_SLOT_LEFT_REAR=3,
               ESPNOW_SLOT_RIGHT_FRONT=4, ESPNOW_SLOT_RIGHT_REAR=5 } espnow_slot_t;
typedef struct __attribute__((packed)) {
    float distance_cm[ESPNOW_SENSOR_SLOT_COUNT];
    uint8_t valid[ESPNOW_SENSOR_SLOT_COUNT];
} espnow_sensor_msg_t;
TBS_STATIC_ASSERT(ESPNOW_SENSOR_SLOT_COUNT == SENSOR_COUNT, ...);
\end{verbatim}

Bản đồ bộ nhớ gói tin 30\,B:

\begin{center}
\small
\renewcommand{\arraystretch}{1.15}
\begin{longtable}{@{}lcc@{}}
\caption{Bố cục payload ESP-NOW.}\\
\toprule
\textbf{Trường} & \textbf{Offset} & \textbf{Kích thước} \\
\midrule
\endhead
\code{distance\_cm[0..5]} & 0 & 6\,$\times$ 4\,B = 24\,B (float) \\
\code{valid[0..5]} & 24 & 6\,$\times$ 1\,B = 6\,B \\
\textbf{Tổng} & 0 & \textbf{30\,B} \\
\bottomrule
\end{longtable}
\end{center}

Lưu ý ánh xạ \textbf{vật lý $\rightarrow$ dây} nằm ở \code{main.cpp:53}:
\code{SENSOR\_ESPNOW\_SLOT[SENSOR\_COUNT]} — \code{SENSOR\_PINS[0]} (FRONT) là
dây slot FRONT, \code{SENSOR\_PINS[5]} (REAR) là dây slot REAR. Nhờ
\code{static\_assert} (R4), số slot trên dây luôn bằng \code{SENSOR\_COUNT}.

\section{espnow\_client.cpp (42 dòng) — phát ESP-NOW}

\begin{itemize}
\item \code{begin()}: \code{WiFi.mode(WIFI\_STA)} + \code{disconnect()};
\code{esp\_wifi\_set\_channel(ESPNOW\_CHANNEL,...)} cố định kênh 1; init
ESP-NOW; đăng ký \code{onDataSent}; thêm peer với MAC đích
(\code{peerInfo.encrypt = false} — không mã hoá, chấp nhận vì mạng cục bộ và
payload không nhạy cảm).
\item \code{sendReading(msg)}: \code{esp\_now\_send(ESPNOW\_PEER\_MAC, ...,
sizeof(msg))} — gửi đúng 30\,B, trả \code{true} nếu đưa vào hàng đợi thành
công.
\end{itemize}

\section{networkTask trong main.cpp — điều phối phát 500\,ms}

Task core 0 (dòng 202--240):
\begin{enumerate}
\item \code{s\_espNowClient.begin()} một lần.
\item Mỗi 500\,ms (\code{now - lastSendMs >= ESPNOW\_SEND\_INTERVAL\_MS}):
tạo \code{espnow\_sensor\_msg\_t msg = \{\}}; với mỗi cảm biến
\code{i = 0..SENSOR\_COUNT-1} đọc \code{sharedStateGet}, nếu hợp lệ thì ghi
\code{msg.distance\_cm[slot]} và \code{msg.valid[slot] = 1} theo
\code{SENSOR\_ESPNOW\_SLOT[i]}; rồi \code{sendReading}.
\item \code{vTaskDelay(20)} giữa các lần kiểm tra — không gửi blocking.
\end{enumerate}

Điểm yếu cố ý: slot không lắp phần cứng giữ \code{valid=0} $\rightarrow$
waveshare hiển thị ``--'' thay vì báo DANGER giả.

\section{espnow\_receiver.c (112 dòng) — nhận và theo dõi liên kết}

\begin{itemize}
\item \code{on\_data\_recv}: \textbf{từ chối gói sai kích thước}
(\code{len != sizeof(espnow\_sensor\_msg\_t)}) — ngăn buffer-overread; copy gói
rồi cập nhật timestamp \code{s\_slot\_rx\_us[i]} cho từng slot có
\code{valid[i]}; gọi callback driver với RSSI.
\item \code{espnow\_receiver\_init(cb)}: đặt kênh, init ESP-NOW (chấp nhận trả
\code{ESP\_ERR\_ESPNOW\_EXIST}), cài callback.
\item \code{espnow\_receiver\_is\_linked()}: so \code{now - s\_last\_any\_rx\_us}
với \code{ESPNOW\_LINK\_TIMEOUT\_MS * 1000} (1500\,ms) — nguồn cho badge
``ESP-NOW: LINKED/NO LINK'' trên UI (xem Chương 5).
\end{itemize}

\section{MQTT — CoreiotClient (sensor-node, Arduino + PubSubClient)}

File \code{plugins/coreiot/coreiot\_client.cpp} — đường phụ chỉ build khi
\code{USE\_COREIOT=1} (env \code{yolo\_uno\_coreiot}, R6):
\begin{itemize}
\item \code{begin()}: \code{setServer(COREIOT\_BROKER, COREIOT\_PORT)};
\code{WiFi.begin(WIFI\_SSID, WIFI\_PASSWORD)} — credential từ
\code{credentials.h} \textbf{sinh tự động, gitignored} (R1).
\item \code{ensureWifiConnected()}: không gọi lại \code{WiFi.begin()} nền
(chỉ tái kết nối nền) để tránh reset trạng thái đang thử.
\item \code{ensureMqttConnected()}: thử lại mỗi
\code{COREIOT\_MQTT\_RETRY\_INTERVAL\_MS = 3000}; kết nối bằng
\code{s\_mqttClient.connect("sensor-node", SENSOR\_NODE\_DEVICE\_TOKEN,
nullptr)} — \textbf{access token là username MQTT}, password rỗng (chuẩn
ThingsBoard/CoreIoT).
\item \code{publishTelemetry(payload)}: publish lên
\code{COREIOT\_TELEMETRY\_TOPIC} (\code{v1/devices/me/telemetry}).
\end{itemize}

\subsection{coreiotTask — payload telemetry 500\,ms (main.cpp:252--290)}

\begin{verbatim}
{"d1":25.0,"d2":120.0,...,"d6":400.0,"nearest_cm":25.0,"has_nearest":true}
\end{verbatim}
Dùng \code{snprintf} JSON đơn giản (không thư viện JSON Arduino), đọc từng slot
qua \code{sharedStateGetValue} + gần nhất qua \code{sharedStateGetNearest};
publish thất bại thì bỏ qua (loop() tự duy trì kết nối) — \textbf{không làm hỏng
route cảnh báo chính}.

\section{MQTT — coreiot\_client (waveshare-screen, ESP-IDF esp-mqtt)}

File \code{components/coreiot\_client/coreiot\_client.c} — client đầy đủ hơn
vì screen vừa gửi vừa nhận:
\begin{itemize}
\item \code{mqtt\_event\_handler} — khi \code{MQTT\_EVENT\_CONNECTED}:
  \begin{itemize}
  \item Subscribe \code{v1/devices/me/telemetry},
\code{v1/devices/me/attributes} và \code{v1/devices/me/attributes/response/+}
(QoS 1) — nhận cảnh báo real-time từ Rule-Chain (đường phụ).
  \item Publish telemetry \code{reboot\_event} (\{status ONLINE, reason,
wifi\_ssid\}) để dashboard biết màn hình vừa khởi động.
  \end{itemize}
\item \code{MQTT\_EVENT\_DATA}: ghi \code{s\_last\_mqtt\_rx\_time\_ms} (dùng cho
badge MQTT trên UI) và gọi \code{s\_data\_cb} cho lớp Ứng dụng.
\item \code{MQTT\_EVENT\_DISCONNECTED}: báo \code{s\_mqtt\_cb(false)}; kết nối
lại do esp-mqtt tự thử.
\end{itemize}

\section{Tập trung: luồng mạng 1 chu kỳ}

\begin{verbatim}
SensorTask (100 ms) -> shared_state (mutex)
   |-- networkTask (500 ms): khoi tao msg 30B -> esp_now_send -> screen
   `-- coreiotTask (500 ms): snprintf JSON d1..d6+nearest
              -> PubSubClient.publish(v1/devices/me/telemetry)
              -> broker app.coreiot.io:1883 -> Lớp 3 (Chương 3)
\end{verbatim}

\section{Kết luận lớp Mạng}

Mạng đóng gói đúng hợp đồng (R2), gửi đều nhịp, không blocking, không phụ
thuộc nhau: mất MQTT không ảnh hưởng ESP-NOW. Trạng thái liên kết (link timeout
1500\,ms) và thời gian MQTT rx được chuyển lên Lớp Ứng dụng (Chương 5) để hiển
thị badge.